\section{Operasi Logika Dasar}
\label{sec:logika-dasar}

Instruksi aritmatika memandang operand sebagai satu bilangan utuh dan
menjalarkan \eng{carry} antarbit. Instruksi logika bekerja berbeda:
setiap posisi bit diperlakukan sebagai satuan mandiri yang tidak
memengaruhi posisi tetangganya. Operasi semacam ini disebut
\term{operasi logika} atau \eng{bitwise operation}.

Bagian ini diringkas dari \cite{gunawan2019}, Bagian 7,
hlm.~68--73, dan \cite{sto2001}, BAB X, hlm.~47--51.

\subsection{Tabel Kebenaran}

Empat instruksi logika 8086 bekerja pada operand 8 bit maupun 16 bit.
Hasilnya hanya bergantung pada pasangan bit yang bersesuaian,
seperti ditunjukkan pada Tabel~\ref{tab:kebenaran-logika}.

\begin{table}[H]
    \centering
    \caption{Tabel kebenaran operasi logika 8086.}
    \label{tab:kebenaran-logika}
    \small
    \begin{tabular}{@{}ccccc@{}}
        \toprule
        \textbf{Bit $A$} & \textbf{Bit $B$} &
        \textbf{\asm{AND}} & \textbf{\asm{OR}} &
        \textbf{\asm{XOR}} \\
        \midrule
        0 & 0 & 0 & 0 & 0 \\
        0 & 1 & 0 & 1 & 1 \\
        1 & 0 & 0 & 1 & 1 \\
        1 & 1 & 1 & 1 & 0 \\
        \bottomrule
    \end{tabular}
\end{table}

\subsection{AND, OR, XOR, dan NOT}

\begin{description}
    \item[\asm{AND}] Menghasilkan 1 hanya bila kedua bit bernilai 1.
          Instruksi ini dipakai untuk mempertahankan bit tertentu dan
          memadamkan bit lainnya.

    \item[\asm{OR}] Menghasilkan 1 bila salah satu atau kedua bit
          bernilai 1. Instruksi ini dipakai untuk menyalakan bit
          tertentu tanpa mengubah bit lain.

    \item[\asm{XOR}] Menghasilkan 1 bila kedua bit berbeda.
          Instruksi ini dipakai untuk membalik bit tertentu.

    \item[\asm{NOT}] Membalik seluruh bit operand. Instruksi ini
          tidak memengaruhi register \eng{flag}.
\end{description}

Contoh:

\begin{lstlisting}[language=MASM, numbers=none]
MOV AL, 10100101b       ; AL = A5h
AND AL, 00001111b       ; AL = 05h (nibble rendah dipertahankan)
MOV BL, 11000011b       ; BL = C3h
OR  AL, BL              ; AL = C7h (bit tertentu dinyalakan)
MOV CL, 10100101b       ; CL = A5h
XOR CL, 11111111b       ; CL = 5Ah (seluruh bit dibalik)
NOT CL                  ; CL = A5h (kembali semula)
\end{lstlisting}

\subsection{Pengaruh terhadap Flag}

Operasi logika tidak menghasilkan \eng{carry} antarbit. Karena itu,
\flag{CF} dan \flag{OF} selalu dikosongkan. Flag lain mengikuti
hasil operasi seperti ditunjukkan pada
Tabel~\ref{tab:flag-logika}.

\begin{table}[H]
    \centering
    \caption{Pengaruh operasi logika terhadap flag.}
    \label{tab:flag-logika}
    \small
    \begin{tabular}{@{}p{1.6cm}p{5.2cm}p{6.4cm}@{}}
        \toprule
        \textbf{Flag} & \textbf{Nilai setelah operasi logika} &
        \textbf{Alasan} \\
        \midrule
        \flag{CF} & Selalu 0 & Tidak ada \eng{carry} yang keluar dari
        bit paling kiri \\
        \addlinespace
        \flag{OF} & Selalu 0 & Operasi bitwise tidak menghasilkan
        luapan bertanda \\
        \addlinespace
        \flag{ZF} & 1 bila seluruh bit hasil bernilai 0 &
        Dipakai untuk memeriksa apakah suatu bidang bit kosong \\
        \addlinespace
        \flag{SF} & Menyalin bit paling kiri hasil &
        Menandakan hasil dianggap negatif bila ditafsirkan bertanda \\
        \addlinespace
        \flag{PF} & 1 bila jumlah bit 1 pada byte rendah genap &
        Dipakai pada pemeriksaan paritas \\
        \bottomrule
    \end{tabular}
\end{table}

\begin{infobox}
\vspace{1.1em}
\asm{NOT} adalah pengecualian: instruksi ini tidak menyentuh
\eng{flag} sama sekali. Bila hasil \asm{NOT} perlu diuji, tambahkan
\lstasm|TEST AL, AL| sesudahnya agar \eng{flag} diperbarui
berdasarkan isi \reg{AL} yang baru.
\end{infobox}

\begin{peringatanbox}
\vspace{1.1em}
Jangan mengandalkan \flag{CF} setelah operasi logika. Berbeda dari
\asm{ADD} atau \asm{SUB}, \asm{AND}, \asm{OR}, dan \asm{XOR} selalu
mengosongkan \flag{CF}. Nilai \flag{CF} yang tersisa dari operasi
sebelumnya tidak dapat dipakai setelah operasi logika.
\end{peringatanbox}
